\documentclass[11pt,a4paper]{article}
\usepackage[margin=2.5cm]{geometry}
\usepackage{amsmath,amssymb}
\usepackage{booktabs}
\usepackage{array}
\usepackage{parskip}
\setlength{\parskip}{8pt}
\usepackage[hidelinks]{hyperref}

\newcommand{\cd}{c^\dagger}
\newcommand{\ad}{a^\dagger}
\newcommand{\up}{\uparrow}
\newcommand{\dn}{\downarrow}

\begin{document}

\begin{center}
{\large VQE for the Hubbard ring on MIMIQ: the equations, step by step}\\[4pt]
Khuram Shahzad, October 2026
\end{center}

The ground state of the Hubbard ring is prepared with the variational quantum
eigensolver (VQE) on MIMIQ and passed on to the $n$BRDM and GME analysis. All
the equations needed for this are given below, step by step. Sites $k$ and
orbitals $p$ are counted from 1, qubits from 0. Every formula has been checked
numerically (Section~\ref{sec:checks}).

\emph{Symbols.} $t, U, V$: model parameters; $\mathcal{U}(\boldsymbol\theta)$:
the UCCSD circuit; $R_P(\varphi)$: the Pauli rotation gate;
$\mathcal{R}_C$: the rotation from the orbitals to the sites; $q$: a qubit
index; $c_{k,\sigma}$ ($a_{p,\sigma}$): annihilation operator on site $k$
(orbital $p$) with spin $\sigma$.

%-----------------------------------------------------------------------------
\section*{The procedure in short}

\emph{One Hamiltonian.} The Hamiltonian is always the full $H_e$ of
Eq.~(\ref{eq:H}), with $t$, $U$ and $V$. It is never changed, and VQE never
diagonalizes it: VQE only measures its expectation value. (The exact
diagonalization, ED, of $H_e$ is used only as the reference to compare with.)

\emph{Two bases for the circuit.} The qubits can stand for two different sets
of one-electron states: the \emph{sites} (an electron on site $k$) or the
\emph{orbitals} (an electron in a wave spread over the ring). Both describe the
same space; a change of basis changes how states look, not the physics. The
orbitals are used only to \emph{build the starting state}, because there it is
a few $X$ gates; the energy and the $n$BRDMs are always measured on the sites.

\emph{The steps, in order.}
\begin{enumerate}
  \item Hamiltonian $H_e(t,U,V)$ on the sites (Section~\ref{sec:H}).
  \item One qubit per spin orbital and the Jordan--Wigner transformation:
    $H_e$ becomes a sum of Pauli strings on the sites (Section~\ref{sec:jw}).
  \item Orbitals from the $L \times L$ one-electron hopping matrix, only to
    prepare the state (Section~\ref{sec:orbitals}).
  \item Starting state: Hartree--Fock, the lowest orbitals filled, i.e.\ $X$
    gates (Section~\ref{sec:hf}).
  \item Ansatz: UCCSD with parameters $\boldsymbol\theta$, acting on the
    orbitals (Section~\ref{sec:uccsd}).
  \item Rotation of the state from the orbitals to the sites
    (Section~\ref{sec:rotation}).
  \item Energy $E(\boldsymbol\theta) = \langle H_e \rangle$ measured on MIMIQ,
    with $U$ and $V$ included (Section~\ref{sec:energy}).
  \item COBYLA changes $\boldsymbol\theta$ until $E(\boldsymbol\theta)$ stops
    decreasing (Section~\ref{sec:opt}).
  \item From the final state: $n$BRDMs and GME, compared with ED
    (Section~\ref{sec:gme}).
\end{enumerate}
The circuit is therefore
\begin{equation*}
  |0\cdots0\rangle
  \xrightarrow{\ X \text{ gates}\ } |\Phi_0\rangle
  \xrightarrow{\ \mathcal{U}(\boldsymbol\theta)\ } |\Psi_{\rm orb}(\boldsymbol\theta)\rangle
  \xrightarrow{\ \mathcal{R}_C\ } |\Psi_e(\boldsymbol\theta)\rangle
  \ \to\ \text{measure } \langle H_e\rangle ,
\end{equation*}
where the first three objects are written in the orbitals and the last one on
the sites.

%-----------------------------------------------------------------------------
\section{The Hubbard Hamiltonian}\label{sec:H}

The one-dimensional extended Hubbard model with $K$ electrons on a ring of $L$
sites is defined by the Hamiltonian
\begin{gather}
{H}_e = - t \sum_{k=1}^L \sum_{\sigma=\uparrow,\downarrow} \big({c}_{k,\sigma}^\dagger\,{c}_{k+1,\sigma} + {c}_{k+1,\sigma}^\dagger\,{c}_{k,\sigma}\big)
+ U \sum_{k=1}^L {n}_{k,\uparrow}\,{n}_{k,\downarrow} \nonumber\\
+ V \sum_{k=1}^L {n}_{k}\,{n}_{k+1} + \epsilon \sum_{k=1}^L {n}_{k}\,.
\label{eq:H}
\end{gather}
Here $n_{k,\sigma} = \cd_{k,\sigma} c_{k,\sigma}$ is the number operator of
site $k$ ($1 \le k \le L$) with spin $\sigma = \up, \dn$, and
$n_k = n_{k,\up} + n_{k,\dn}$ the number of electrons on site $k$; $t$, $U$ and
$V$ are the hopping, on-site and nearest-neighbor interaction parameters. The
sum over sites is written out for each term: the hopping has two sums, over the bonds $(k, k+1)$ and over the
spin, since an electron of either spin can hop and the hopping does not flip
the spin. Because of the periodic boundary conditions, site $L+1$ coincides
with site $1$. The detection orbital $\phi_k$ is the one localized on site $k$,
and the spins are extracted from $n$ consecutive orbitals
$\phi_1, \dots, \phi_n$. At fixed $K$ the $\epsilon$ term is the constant
$\epsilon K$; we use $\epsilon = 0$. We work
at half filling, $K = L$, in the sector with $K_\up = K_\dn = K/2$ electrons of
each spin.

Taking $t$ out of every term gives $H_e(t,U,V) = t\, H_e(1, U/t, V/t)$.
Multiplying by $t > 0$ rescales all energies by $t$ but leaves the eigenstates
unchanged. Hence the ground state $|\Psi_0\rangle$ depends only on $U/t$ and
$V/t$, and its energy is $E_0(t,U,V) = t\, E_0(1, U/t, V/t)$.

%-----------------------------------------------------------------------------
\section{Spin orbitals, qubits and Jordan--Wigner}\label{sec:jw}

Each spin orbital $(k,\sigma)$ is stored on one qubit; qubits are counted
from 0:
\begin{equation}\label{eq:qubit}
  q(k,\sigma) = 2(k-1) + \sigma , \qquad \sigma = 0 \ (\up), \ 1 \ (\dn),
\end{equation}
so that site $k$ is stored on the two qubits with labels $2k-2$ ($\up$) and
$2k-1$ ($\dn$). The labels run from $0$ to $2L-1$, so the ring needs $2L$
qubits (12 for $L = 6$, labels $0,\dots,11$). A qubit in
$|1\rangle$ means the spin orbital is occupied. The Jordan--Wigner
transformation writes the operators of Eq.~(\ref{eq:H}) with Pauli matrices on
these qubits. With $q = q(k,\sigma)$,
\begin{equation}\label{eq:jw}
  c_{k,\sigma} = Z_0 Z_1 \cdots Z_{q-1}\, \frac{X_q + i Y_q}{2} , \qquad
  n_{k,\sigma} = \cd_{k,\sigma} c_{k,\sigma} = \frac{1 - Z_q}{2} .
\end{equation}
Here $c_{k,\sigma}$ removes the electron from the spin orbital $(k,\sigma)$
(it maps $|1\rangle_q \to |0\rangle_q$), and $n_{k,\sigma}$ is its occupation,
$0$ (empty, $Z_q = +1$) or $1$ (occupied, $Z_q = -1$); the string
$Z_0 \cdots Z_{q-1}$ gives the fermion sign. For the three types of term in
Eq.~(\ref{eq:H}) this gives, with $q = q(k,\sigma)$ and $q' = q(k+1,\sigma)$
in the hopping, ordered so that $q < q'$,
\begin{align}
  \cd_{k,\sigma} c_{k+1,\sigma} + \cd_{k+1,\sigma} c_{k,\sigma}
     &= \tfrac12 \left( X_q Z_{q+1} \cdots Z_{q'-1} X_{q'}
                      + Y_q Z_{q+1} \cdots Z_{q'-1} Y_{q'} \right) , \label{eq:jw-hop}\\
  n_{k,\sigma}\, n_{k',\sigma'}
     &= \tfrac14 \left( 1 - Z_q - Z_{q'} + Z_q Z_{q'} \right) ,
        \qquad q = q(k,\sigma) \ne q' = q(k',\sigma') , \label{eq:jw-nn}
\end{align}
where Eq.~(\ref{eq:jw-nn}) gives the $U$ term with $k' = k$, $\sigma = \up$,
$\sigma' = \dn$, and the $V$ term through
$n_k n_{k+1} = \sum_{\sigma,\sigma'} n_{k,\sigma}\, n_{k+1,\sigma'}$.
In Eq.~(\ref{eq:jw-hop}), for $k < L$, $q' = q + 2$ and the string is the single
$Z_{q+1}$ of the other spin on site $k$; for the ring bond ($k = L$,
$k+1 \equiv 1$) the two qubits are $\sigma$ and $2L-2+\sigma$, and the string
covers all the qubits in between. Example: hopping of $\up$ between sites 1
and 2 is $\tfrac12 (X_0 Z_1 X_2 + Y_0 Z_1 Y_2)$.
The Hamiltonian therefore becomes a sum of Pauli strings with real coefficients,
\begin{equation}\label{eq:H-pauli}
  H_e = \lambda_0 + \sum_{\mu} \lambda_\mu P_\mu , \qquad P_\mu \in \{I, X, Y, Z\}^{\otimes 2L} ,
\end{equation}
with the constant $\lambda_0 = UL/4 + VL$ at $\epsilon = 0$.
For $L = 6$ there are 42 strings at $(t,U,V) = (1,4,0)$ (24 hopping, 18 on-site)
and 66 at $(1,10,5)$.

Bit order: MIMIQ takes qubit 0 as the lowest bit of a basis state. This
matters only when a state vector is loaded or read out and compared with
another convention; gates and Pauli strings refer to qubit indices and do not
depend on it.

%-----------------------------------------------------------------------------
\section{Orbitals for preparing the state}\label{sec:orbitals}

This section does not change the Hamiltonian: $H_e$ stays the full
Hamiltonian of Eq.~(\ref{eq:H}), with $U$ and $V$. What is needed here is a
starting state for the circuit that is easy to prepare with gates and already
close to the ground state. It is built from one-electron states (orbitals)
obtained from the hopping term alone, for three reasons:
\begin{itemize}
  \item the hopping is the only one-electron part of $H_e$, so it defines
    one-electron states; $U$ and $V$ are interactions between electrons, and
    they are already diagonal on the sites (they only count electrons);
  \item at $U = V = 0$, $H_e$ is the hopping alone, and its exact ground state
    is obtained by filling the lowest orbitals (Section~\ref{sec:hf}); for
    $U \ne 0$ this state is the starting point that VQE improves;
  \item finding the orbitals needs only the diagonalization of a small
    $L \times L$ matrix, whereas the exact ground state needs the
    diagonalization of $H_e$ in the full many-electron space (924 states for
    $L = 6$, growing exponentially with $L$), which VQE avoids.
\end{itemize}
The hopping matrix of the ring (one electron, $L \times L$) is
\begin{equation}
  h_{kk'} = -t \left( \delta_{k',k+1} + \delta_{k',k-1} \right) , \qquad
  1 \le k, k' \le L \ \ (\text{site } L+1 \equiv 1),
\end{equation}
and its eigenvectors define the orbitals $p = 1, \dots, L$, sorted by energy,
\begin{equation}\label{eq:orbitals}
  h = C \, \mathrm{diag}(e_1, \dots, e_L) \, C^{T} , \qquad
  c_{k,\sigma} = \sum_{p=1}^{L} C_{kp}\, a_{p,\sigma} ,
\end{equation}
with $C$ real and orthogonal (row $k$: site, column $p$: orbital). The
columns of $C$ are the eigenvectors of $h$; they form the new basis (the
orbitals), and the eigenvalues $e_p$ are the orbital energies. For $L = 6$:
$e_p = -2t, -t, -t, t, t, 2t$; the lowest orbital is spread evenly over the
ring, $C_{k1} = 1/\sqrt{6}$, and the higher ones change sign more often from
site to site. The orbitals and the sites are two bases of the same
one-electron space, related by the orthogonal matrix $C$; nothing is lost in
the change. The orbitals depend only on $t$, not on $U$ or $V$; they play the
role of the Hartree--Fock orbitals of a molecule. In the
orbital basis the spin orbital $(p,\sigma)$ is stored on qubit
$q(p,\sigma) = 2(p-1) + \sigma$, the same rule as Eq.~(\ref{eq:qubit}).

In the orbital basis the hopping is diagonal but the $U$ and $V$ terms are
not: writing $H_e$ in the orbitals (substituting Eq.~(\ref{eq:orbitals}) into
Eq.~(\ref{eq:H})) gives the same spectrum but many more Pauli strings
(306 instead of 42 at $(1,4,0)$). On the sites it is the other way round. No
basis makes both parts diagonal; this is why the problem is hard. The state is
therefore prepared in the orbitals and rotated back to the sites
(Section~\ref{sec:rotation}) before the full $H_e$ is measured.

%-----------------------------------------------------------------------------
\section{Hartree--Fock reference}\label{sec:hf}

The reference state fills the $K/2$ lowest orbitals with both spins,
\begin{equation}\label{eq:hf}
  |\Phi_0\rangle = \prod_{p=1}^{K/2} \ad_{p,\up}\, \ad_{p,\dn} \, |\mathrm{vac}\rangle ,
\end{equation}
that is, the qubits $0, \dots, K-1$ are set to $|1\rangle$ by $X$ gates. Since
the density is uniform ($\langle n_{k,\sigma}\rangle = 1/2$) and the on-site term
has no exchange part,
\begin{equation}\label{eq:ehf}
  E_{\rm HF} = 2 \sum_{p=1}^{K/2} e_p + U \sum_{k=1}^{L} \langle n_{k,\up}\rangle\langle n_{k,\dn}\rangle
            = 2 \sum_{p=1}^{K/2} e_p + \frac{UL}{4}
            \;\stackrel{L=6}{=}\; -8t + \tfrac32 U \qquad (\epsilon = V = 0).
\end{equation}
$|\Phi_0\rangle$ is the exact ground state of $H_e$ only at $U = V = 0$; for
$U > 0$ it is the starting point of VQE, and Eq.~(\ref{eq:ehf}) is its energy
with the full $H_e$. The reference must be a closed shell: the last filled level must not be
degenerate with the first empty one. This holds for $L = 6$ (levels
$-2t, -t, -t \,|\, t, t, 2t$) but not for $L = 4$ ($-2t, 0 \,|\, 0, 2t$).
At $U/t = 1, 4, 10$, $|\Phi_0\rangle$ (rotated to the sites) has the fidelity
$0.980$, $0.689$ and $0.265$ with the exact ground state, against 3--19\,\% for
the best single configuration on sites.

%-----------------------------------------------------------------------------
\section{The ansatz: UCCSD in the orbital basis}\label{sec:uccsd}

Let $i,j$ run over the occupied orbitals ($1, \dots, K/2$) and $a,b$ over the
empty ones ($K/2+1, \dots, L$). The cluster operator contains all the
spin-conserving single and double excitations,
\begin{equation}\label{eq:T}
  T(\boldsymbol\theta) = \sum_{\sigma} \sum_{a,i} \theta^{\sigma}_{ai}\, \ad_{a,\sigma} a_{i,\sigma}
  + \sum_{\text{doubles}} \theta_{a\sigma\, i\sigma\, b\sigma'\, j\sigma'}\;
    \ad_{a,\sigma} a_{i,\sigma}\, \ad_{b,\sigma'} a_{j,\sigma'} ,
\end{equation}
where the doubles are of three kinds: $\up\dn$ (all $a,i,b,j$), $\up\up$ and
$\dn\dn$ ($a<b$, $i<j$). With $o = K/2$ occupied and $v = L - K/2$ empty orbitals
the number of parameters is
\begin{equation}\label{eq:npar}
  N_\theta = \underbrace{2ov}_{\text{singles}}
           + \underbrace{(ov)^2}_{\up\dn}
           + \underbrace{2\tbinom{o}{2}\tbinom{v}{2}}_{\up\up,\ \dn\dn}
  \;\stackrel{o=v=3}{=}\; 18 + 81 + 18 = 117 .
\end{equation}
The parameters are ordered as singles $\up$, singles $\dn$, doubles $\up\dn$,
$\up\up$, $\dn\dn$. The UCCSD state is
\begin{equation}\label{eq:uccsd}
  |\Psi_{\rm orb}(\boldsymbol\theta)\rangle
  = e^{\,T(\boldsymbol\theta) - T(\boldsymbol\theta)^\dagger} |\Phi_0\rangle .
\end{equation}
Since $T - T^\dagger$ is anti-Hermitian, $G = i(T - T^\dagger)$ is Hermitian, and
Jordan--Wigner turns it into Pauli strings with real coefficients that are
linear in $\boldsymbol\theta$:
\begin{equation}\label{eq:G}
  e^{T - T^\dagger} = e^{-iG} , \qquad
  G(\boldsymbol\theta) = \sum_m \gamma_m(\boldsymbol\theta)\, Q_m .
\end{equation}
On MIMIQ the exponential is applied with one first-order Trotter step,
\begin{equation}\label{eq:trotter}
  \mathcal{U}(\boldsymbol\theta) = \prod_m e^{-i \gamma_m Q_m} = \prod_m R_{Q_m}(2 \gamma_m) ,
  \qquad R_P(\varphi) = e^{-i\varphi P/2} ,
\end{equation}
where $R_P$ is the Pauli rotation gate. The strings $Q_m$ do not all
commute, so Eq.~(\ref{eq:trotter}) is not exactly Eq.~(\ref{eq:uccsd}); it
\emph{defines} the ansatz, and the optimizer adjusts
$\boldsymbol\theta$ for it. Every excitation conserves the number of electrons
and $S_z$, so every factor of Eq.~(\ref{eq:trotter}) does too; the total spin
$S^2$ is not conserved, because the 117 amplitudes are independent (not spin
adapted). The circuit starts with the $X$ gates of Eq.~(\ref{eq:hf}). After decomposition into
one- and two-qubit gates the circuit has 29\,298 gates (9\,024 CNOTs) for
$L = 6$.

%-----------------------------------------------------------------------------
\section{Rotation from the orbitals back to the sites}\label{sec:rotation}

$|\Psi_{\rm orb}\rangle$ is written in the orbital modes $a_{p,\sigma}$. We need
the same physical state written in the site modes $c_{k,\sigma}$, on the same
qubits. This is done by the unitary $\mathcal{R}_C$ defined by
\begin{equation}\label{eq:rc}
  \mathcal{R}_C\, \ad_{p,\sigma}\, \mathcal{R}_C^\dagger = \sum_{k=1}^{L} C_{kp}\, \cd_{k,\sigma} , \qquad
  \mathcal{R}_C |\mathrm{vac}\rangle = |\mathrm{vac}\rangle ,
\end{equation}
which replaces every orbital creation operator by the combination of site
operators of Eq.~(\ref{eq:orbitals}). The state on the sites is then
\begin{equation}\label{eq:psi}
  |\Psi_e(\boldsymbol\theta)\rangle = \mathcal{R}_C\, \mathcal{U}(\boldsymbol\theta)\, X_0 X_1 \cdots X_{K-1} |0\cdots0\rangle .
\end{equation}

\emph{Givens decomposition.} The orthogonal matrix $C$ is brought to a diagonal
matrix of signs by rotations $W_r$ of two neighboring rows $(p, p+1)$,
\begin{equation}\label{eq:givens}
  W_{N_g} \cdots W_2\, W_1\, C = D , \qquad D = \mathrm{diag}(\pm 1) ,
  \qquad N_g = \tfrac{L(L-1)}{2} = 15 ,
\end{equation}
each rotation zeroing one element below the diagonal.
A rotation by $\vartheta$ of the orbitals $p$ and $p+1$ is, for both spins, the
fermionic unitary
\begin{equation}
  g_p(\vartheta) = \exp\!\Big[ \vartheta \sum_\sigma
     \big( \ad_{p+1,\sigma} a_{p,\sigma} - \ad_{p,\sigma} a_{p+1,\sigma} \big) \Big] .
\end{equation}
With $q = q(p,\sigma)$, the two orbitals sit on the qubits $q$ and $q+2$, with
the other spin on $q+1$ in between, and Jordan--Wigner gives exactly two
commuting strings,
\begin{equation}\label{eq:givens-jw}
  i\vartheta \big( \ad_{p+1,\sigma} a_{p,\sigma} - \ad_{p,\sigma} a_{p+1,\sigma} \big)
  = \frac{\vartheta}{2} \big( X_q Z_{q+1} Y_{q+2} - Y_q Z_{q+1} X_{q+2} \big) ,
\end{equation}
so that, without any Trotter error,
\begin{equation}
  g_p(\vartheta) = \prod_{\sigma}
    R_{X_q Z_{q+1} Y_{q+2}}(\vartheta)\, R_{Y_q Z_{q+1} X_{q+2}}(-\vartheta) .
\end{equation}
The whole rotation $\mathcal{R}_C$ is a $Z$ gate on the qubits $q(p,\up)$ and
$q(p,\dn)$ for every $D_{pp} = -1$, followed by the $N_g$ rotations of
Eq.~(\ref{eq:givens}) in reverse order: $4 N_g = 60$ Pauli rotations for
$L = 6$. Which signs occur depends on the
numerical eigenvectors: at $t = 1$ all $D_{pp} = +1$, at $t = 0.05$
$D_{66} = -1$.

Why rotate instead of measuring in the orbital basis: the energy
then needs only the site strings of Eq.~(\ref{eq:H-pauli}) (42 instead of 306),
the $n$BRDMs are defined on sites and can be measured directly, and the
Hamiltonian is used unchanged.

%-----------------------------------------------------------------------------
\section{The energy measured on MIMIQ}\label{sec:energy}

For given $\boldsymbol\theta$ the circuit of Eq.~(\ref{eq:psi}) is run on
MIMIQ, and the only output is the expectation value
\begin{equation}\label{eq:energy}
  E(\boldsymbol\theta) = \langle \Psi_e(\boldsymbol\theta) | H_e | \Psi_e(\boldsymbol\theta) \rangle
  = \lambda_0 + \sum_\mu \lambda_\mu \langle P_\mu \rangle_{\boldsymbol\theta} ,
\end{equation}
obtained with one expectation-value instruction for the Hamiltonian of
Eq.~(\ref{eq:H-pauli}), the constant $\lambda_0$ added
separately. On the exact simulator (Exaqt) $\langle P_\mu\rangle$ is exact; on
the digital twin or a device each $\langle P_\mu\rangle$ is estimated from a
finite number of shots. By the variational principle
$E(\boldsymbol\theta) \ge E_0$ for every $\boldsymbol\theta$.

%-----------------------------------------------------------------------------
\section{The optimization}\label{sec:opt}

The parameters are found by minimizing Eq.~(\ref{eq:energy}),
\begin{equation}
  \boldsymbol\theta^* = \arg\min_{\boldsymbol\theta} E(\boldsymbol\theta) ,
  \qquad E_{\rm VQE} = E(\boldsymbol\theta^*) ,
\end{equation}
with COBYLA, a derivative-free optimizer that needs only the energy values:
\begin{itemize}
  \item start: $\boldsymbol\theta = 0$, that is $|\Phi_0\rangle$ and
    $E = E_{\rm HF}$;
  \item chunks of 600 COBYLA iterations ($\rho_{\rm beg} = 0.2$, tolerance
    $10^{-10}$); each chunk restarts from the best point so far, moved by a
    random step of size $5\times10^{-4}$;
  \item stop when the best energy has improved by less than $10^{-6}$ in 3
    consecutive chunks, or after at most 25 chunks (15\,000 energy evaluations).
\end{itemize}
CCSD amplitudes are not used as a start: with $t = 1$ CCSD gives
$E_{\rm CCSD} = -6.6011$, $-3.7171$ and $-5.5382$ at $U = 1$, $4$ and $10$,
against $E_0 = -6.6012$, $-3.6687$ and $-1.6644$; at large $U$ it is far below
$E_0$ (it is not variational).

%-----------------------------------------------------------------------------
\section{After convergence: $n$BRDMs and GME}\label{sec:gme}

The state $|\Psi_e(\boldsymbol\theta^*)\rangle$ is prepared again with the same
circuit, and the $n$BRDM of each cluster of $n$ consecutive sites is measured
(orbitals $\phi_1, \dots, \phi_n$ and their translations) from Pauli strings, exactly as for the exact state:
\begin{equation}
  \Gamma^{(n)}_{\boldsymbol\alpha,\boldsymbol\alpha'}
  = \sum_\nu w^{\boldsymbol\alpha\boldsymbol\alpha'}_\nu \langle P_\nu \rangle_{\boldsymbol\theta^*} ,
  \qquad \rho_s = \frac{\Gamma^{(n)}}{\mathrm{Tr}\,\Gamma^{(n)}} ,
\end{equation}
where $\boldsymbol\alpha, \boldsymbol\alpha'$ label the spin configurations of
the $n$ sites, $P_\nu$ are the Pauli strings on the cluster's $2n$ qubits and
$w_\nu$ their weights. $\mathcal{C}_{\rm GME}(\rho_s)$ is computed with the GCG
convex-roof solver. Because the VQE state may break the translation symmetry
slightly, all $L$ clusters are measured for each $n$. For every parameter point
we report side by side
\begin{equation}
  \Delta E = E_{\rm VQE} - E_0 , \quad
  F = |\langle \Psi_0 | \Psi_e(\boldsymbol\theta^*) \rangle|^2 , \quad
  \max |\Gamma^{(n)}_{\rm VQE} - \Gamma^{(n)}_{\rm ED}| , \quad
  |\mathcal{C}_{\rm GME}^{\rm VQE} - \mathcal{C}_{\rm GME}^{\rm ED}| ,
\end{equation}
where the fidelity $F$ uses the state vector and is available only on the
simulator. The parameter points are the $U$-scan
($t = 1$, $V = 0$, $U = 1, \dots, 20$), the $t$-scan ($U = 10$, $V = 0$,
$t = 0.05, \dots, 2$) and the $V$-scan ($t = 1$, $U = 10$, $V = 0, \dots, 12$),
65 points in total, each compared with exact diagonalization (ED).

%-----------------------------------------------------------------------------
\section{What has been checked}\label{sec:checks}

\begin{table}[!h]
\centering
\small
\begin{tabular}{>{\raggedright\arraybackslash}p{3.5cm}>{\raggedright\arraybackslash}p{6.3cm}>{\raggedright\arraybackslash}p{4.3cm}}
  \toprule
  Step & Check & Result \\
  \midrule
  $H_e$, Eq.~(\ref{eq:H}) & equal to the formula written term by term; Hermitian; conserves $K$, $S_z$, $S^2$; translation invariant & exact (5 points) \\
  $H_e$ on MIMIQ, Eq.~(\ref{eq:H-pauli}) & $\langle H_e\rangle$ in the exact ground state $= E_0$ & $3\times10^{-15}$ ($L = 4, 6$) \\
  Eqs.~(\ref{eq:jw-hop}), (\ref{eq:jw-nn}) & coefficients one by one; $\langle H_e\rangle$ on a random circuit & $10^{-12}$, $10^{-10}$ \\
  Orbitals, Eq.~(\ref{eq:orbitals}) & $E_0$ equal in both bases & $10^{-10}$ (5 points) \\
  Hartree--Fock, Eq.~(\ref{eq:ehf}) & $E(\boldsymbol\theta = 0) = -8t + \tfrac32 U$ on MIMIQ & $6.6\times10^{-14}$ \\
  Rotation, Eq.~(\ref{eq:rc}) & exact orbital ground state $\to$ site ground state; one electron in orbital $p$ $\to$ amplitudes $C_{kp}$ & $F = 1.000000000000$ (5 points); $10^{-10}$ \\
  Eq.~(\ref{eq:givens-jw}) & Jordan--Wigner of one Givens generator & two strings, $\pm\vartheta/2$ \\
  $N_\theta$, Eq.~(\ref{eq:npar}) & excitations counted one by one & 18 + 99 = 117 \\
  UCCSD, Eq.~(\ref{eq:trotter}) & conserves $K$, $S_z$; MIMIQ energy $=$ state-vector energy $\ge E_0$ & $10^{-10}$ \\
  $n$BRDMs & from Pauli strings on MIMIQ vs.\ ED & $10^{-15}$ ($n \le 4$) \\
  \bottomrule
\end{tabular}
\caption{Checks of the steps above. The five points are
$(t,U,V) = (1,1,0)$, $(1,4,0)$, $(1,10,0)$, $(0.05,10,0)$, $(1,10,5)$.}
\end{table}

\emph{Converged UCCSD runs on MIMIQ.} $L = 6$, $t = 1$, $V = 0$, the full
circuit of Eq.~(\ref{eq:psi}) on Exaqt, COBYLA from $\boldsymbol\theta = 0$.
They agree with UCCSD with the exact exponential of Eq.~(\ref{eq:uccsd})
(no Trotter step, converged with BFGS):
\begin{center}
\begin{tabular}{cccccc}
  \toprule
  & & \multicolumn{2}{c}{exact exponential} & \multicolumn{2}{c}{Trotter step, MIMIQ} \\
  $U$ ($t=1$) & $E_0$ & $\Delta E$ & $F$ & $\Delta E$ & $F$ \\
  \midrule
  1  & $-6.6012$ & $1.0\times10^{-4}$ & 0.99999 & $1.6\times10^{-4}$ & 0.99997 \\
  4  & $-3.6687$ & $3.0\times10^{-2}$ & 0.9937  & -- & -- \\
  10 & $-1.6644$ & $3.9\times10^{-1}$ & 0.943   & $3.9\times10^{-1}$ & 0.942 \\
  \bottomrule
\end{tabular}
\end{center}
UCCSD is accurate at small $U/t$ and degrades at large $U/t$, where the ground
state is far from a single configuration. For large $U/t$ and for the digital
twin (the UCCSD circuit is too deep), shallower ans\"atze such as the
Hamiltonian variational ansatz (HVA) and the local unitary cluster Jastrow
ansatz (LUCJ) are better suited.

%-----------------------------------------------------------------------------
\section{Practical points on MIMIQ}

\begin{itemize}
  \item The exact simulator (Exaqt) does not execute Pauli rotation gates
    directly, also when they sit inside the Trotter block of
    Eq.~(\ref{eq:trotter}); the circuit is first decomposed into standard
    one- and two-qubit gates.
  \item A Pauli rotation needs a real angle; coefficients that come out of the
    Jordan--Wigner transformation as complex numbers with zero imaginary part
    are passed as their real part.
  \item The sign step of $\mathcal{R}_C$ applies $Z$ only where $D_{pp} = -1$;
    it matters, since $D_{pp} = -1$ does occur (e.g.\ at $t = 0.05$).
  \item The cost of one UCCSD energy evaluation for $L = 6$ on Exaqt is about
    0.6~s on one core; a converged run needs $10^3$--$10^4$ evaluations.
\end{itemize}

\end{document}
